\subsection{Основная лемма о линейной зависимости}

\begin{definition}
Пусть $a_1, a_2, \dots, a_k$ — набор векторов из линейного пространства $V$ над полем $F$. \textbf{Линейной комбинацией} этих векторов называется выражение вида:
\[ \lambda_1 a_1 + \lambda_2 a_2 + \dots + \lambda_k a_k \]
где $\lambda_i \in F$ — скаляры. Линейная комбинация называется \textbf{тривиальной}, если все её коэффициенты равны нулю ($\lambda_i = 0$). Если хотя бы один коэффициент отличен от нуля, комбинация называется \textbf{нетривиальной}.
\end{definition}

\begin{definition}
Система векторов $\{a_1, \dots, a_k\}$ называется \textbf{линейно независимой} (ЛНЗ), если их линейная комбинация равна нулевому вектору $\mathbf{0}$ только в тривиальном случае (когда все $\lambda_i = 0$). 
В противном случае (если существует нетривиальный набор коэффициентов, дающий $\mathbf{0}$) система называется \textbf{линейно зависимой} (ЛЗ).
\end{definition}

\begin{theorem}[Основная лемма о линейной зависимости]
Пусть в линейном пространстве $V$ заданы две системы векторов:
\begin{enumerate}
    \item $S_1 = \{a_1, a_2, \dots, a_r\}$ — система из $r$ векторов.
    \item $S_2 = \{b_1, b_2, \dots, b_k\}$ — система из $k$ векторов.
\end{enumerate}
Если система $S_1$ является линейно независимой и каждый вектор из $S_1$ линейно выражается через векторы системы $S_2$, то количество векторов в первой системе не может превышать количества векторов во второй:
\[ r \le k \]
\end{theorem}

\begin{tcolorbox}[title=Доказательство, breakable]
Докажем от противного. Предположим, что $r > k$.
Согласно условию, каждый вектор $a_i$ выражается через векторы $b_j$. Запишем эти разложения:
\[
\begin{cases}
a_1 = \alpha_{11} b_1 + \alpha_{21} b_2 + \dots + \alpha_{k1} b_k \\
a_2 = \alpha_{12} b_1 + \alpha_{22} b_2 + \dots + \alpha_{k2} b_k \\
\dots \\
a_r = \alpha_{1r} b_1 + \alpha_{2r} b_2 + \dots + \alpha_{kr} b_k
\end{cases}
\]
Попробуем построить нетривиальную линейную комбинацию векторов $a_1, \dots, a_r$, равную $\mathbf{0}$. Рассмотрим уравнение:
\[ \lambda_1 a_1 + \lambda_2 a_2 + \dots + \lambda_r a_r = \mathbf{0} \]
Подставим в это равенство выражения для $a_i$ через векторы системы $S_2$:
\[ \lambda_1 \left( \sum_{j=1}^k \alpha_{j1} b_j \right) + \lambda_2 \left( \sum_{j=1}^k \alpha_{j2} b_j \right) + \dots + \lambda_r \left( \sum_{j=1}^k \alpha_{jr} b_j \right) = \mathbf{0} \]
Перегруппируем слагаемые, собирая коэффициенты при каждом векторе $b_j$:
\[ \left( \sum_{i=1}^r \alpha_{1i} \lambda_i \right) b_1 + \left( \sum_{i=1}^r \alpha_{2i} \lambda_i \right) b_2 + \dots + \left( \sum_{i=1}^r \alpha_{ki} \lambda_i \right) b_k = \mathbf{0} \]
Для того чтобы это равенство выполнялось, достаточно потребовать, чтобы коэффициенты при каждом $b_j$ обратились в ноль. Это приводит нас к однородной СЛАУ относительно неизвестных $\lambda_1, \lambda_2, \dots, \lambda_r$:
\[
\begin{cases}
\alpha_{11} \lambda_1 + \alpha_{12} \lambda_2 + \dots + \alpha_{1r} \lambda_r = 0 \\
\alpha_{21} \lambda_1 + \alpha_{22} \lambda_2 + \dots + \alpha_{2r} \lambda_r = 0 \\
\dots \\
\alpha_{k1} \lambda_1 + \alpha_{k2} \lambda_2 + \dots + \alpha_{kr} \lambda_r = 0
\end{cases}
\]
В этой системе $k$ уравнений и $r$ неизвестных. По нашему предположению $r > k$, то есть число неизвестных больше числа уравнений. Согласно следствию из метода Гаусса, такая однородная система обязательно имеет нетривиальное решение $(\lambda^*_1, \dots, \lambda^*_r)$, где не все $\lambda^*_i$ равны нулю.

Следовательно, мы нашли нетривиальную линейную комбинацию векторов $a_1, \dots, a_r$, равную $\mathbf{0}$, что означает линейную зависимость системы $S_1$. Это противоречит условию теоремы о том, что $S_1$ линейно независима. Полученное противоречие доказывает, что наше предположение $r > k$ ложно, и $r \le k$.
\end{tcolorbox}

\begin{tcolorbox}[title=Источники, colback=gray!10]
\begin{itemize}
  \item Лекция №\,3, тема «Основная лемма линейной зависимости», временной отрезок 110:00--115:00
  \item Лекция №\,5, тема «Основная лемма в общем виде», временной отрезок 171:00--173:00
  \item PDF «Линейная алгебра\_compressed (1).pdf», раздел 4, теорема 3 «Основная лемма о линейной зависимости», стр. 15--16 (702--703)
  \item PDF «Вопросы», вопрос №\,7
\end{itemize}
\end{tcolorbox}
